\documentclass[11pt]{article}
\usepackage[polish]{babel}
\usepackage[T1]{fontenc}
\usepackage[utf8]{inputenc}
\usepackage{geometry}
\usepackage{graphicx}
\usepackage{booktabs}
\usepackage{amsmath}
\usepackage{tikz}
\usepackage{caption}
\usepackage{subcaption}
\usepackage{multirow}
\geometry{margin=2.5cm}

% --- biblioteki TikZ do rysunkow logiczno-czasowych (rysunki inline, bez plikow PNG) ---
\usetikzlibrary{shapes.gates.logic.US}
\usetikzlibrary{arrows.meta}
\usetikzlibrary{positioning}
\usetikzlibrary{calc}

\title{Projekt układu sumy liczby 2-bitowej z 3-bitową na bramkach NOR}
\author{}
\date{}

\begin{document}

\maketitle

\section{Opis techniczny projektu}

\subsection{Założenia projektowe}

Celem projektu jest zaprojektowanie układu kombinacyjnego realizującego funkcję sumy dwóch liczb dwubitowych, a następnie zminimalizowanie liczby wykorzystanych bramek logicznych przy zachowaniu realizacji wyłącznie na bramkach NOR.

Projektowany układ będzie dodawał dwie liczby binarne o długości dwóch bitów:
\begin{center}
$A = (A_1 A_0),\quad B = (B_1 B_0)$
\end{center}
Układ posiada cztery wejścia binarne oraz wyjście 3-bitowe. Wynik będzie informował o sumie liczb wraz z ewentualnym przeniesieniem najstarszego bitu.

Przyjęto następujące założenia:
\begin{itemize}
    \item liczby wejściowe mają długość 2 bitów,
    \item wszystkie wejścia przyjmują wartości logiczne 0 lub 1,
    \item układ jest układem kombinacyjnym,
    \item wynik sumy ma 3 bity: \(S_2 S_1 S_0\),
    \item stan wyjścia \(S_2\) oznacza przeniesienie z poziomu bitów starszych,
    \item podstawowym celem projektu jest minimalizacja liczby bramek logicznych,
    \item poprawność działania zostanie sprawdzona za pomocą tablicy prawdy i symulacji.
\end{itemize}

Minimalizacja układu ma na celu zmniejszenie jego złożoności. Mniejsza liczba zastosowanych bramek oznacza prostszy schemat logiczny, mniejsze zapotrzebowanie na elementy oraz potencjalnie mniejsze opóźnienia propagacji sygnału.

\subsection{Specyfikacja wejść i wyjść}

Projektowany sumator posiada cztery wejścia:
\begin{itemize}
    \item \(A_1\) -- starszy bit pierwszej liczby,
    \item \(A_0\) -- młodszy bit pierwszej liczby,
    \item \(B_1\) -- starszy bit drugiej liczby,
    \item \(B_0\) -- młodszy bit drugiej liczby.
\end{itemize}

Wyjścia układu to:
\begin{itemize}
    \item \(S_0\) -- bit wyniku najmłodszy,
    \item \(S_1\) -- bit wyniku pośredni,
    \item \(S_2\) -- najstarszy bit wyniku (przeniesienie).
\end{itemize}

Zależność pomiędzy wejściami a wyjściami można zapisać jako:
\begin{equation}
(A_1 A_0)_2 + (B_1 B_0)_2 = (S_2 S_1 S_0)_2
\end{equation}

Przykładowo, dla:
\begin{center}
$A = 10_2$, $B = 01_2$
\end{center}
otrzymujemy:
\begin{center}
$S = 11_2$
\end{center}

Natomiast dla:
\begin{center}
$A = 11_2$, $B = 11_2$
\end{center}
otrzymujemy:
\begin{center}
$S = 110_2$
\end{center}

Tabela sygnałowa przedstawia podstawowe sygnały wejściowe i wyjściowe projektowanego układu:
\begin{table}[h]
\centering
\caption{Sygnały wejściowe i wyjściowe sumatora 2-bitowego z sumą 3-bitową}
\begin{tabular}{lll}
\toprule
Sygnał & Rodzaj & Opis \\
\midrule
\(A_1\) & wejście & starszy bit pierwszej liczby \\
\(A_0\) & wejście & młodszy bit pierwszej liczby \\
\(B_1\) & wejście & starszy bit drugiej liczby \\
\(B_0\) & wejście & młodszy bit drugiej liczby \\
\(S_2\) & wyjście & najstarszy bit wyniku (przeniesienie) \\
\(S_1\) & wyjście & bit wyniku pośredni \\
\(S_0\) & wyjście & bit wyniku najmłodszy \\
\bottomrule
\end{tabular}
\end{table}

\subsection{Analiza logiczna układu}

Aby wyznaczyć wynik sumy, należy wykonać dodawanie pozycyjne, począwszy od poziomu najmłodszego bitu.

Dla bitów najmłodszych:
\begin{equation}
S_0 = A_0 \oplus B_0
\end{equation}
przy czym w układzie opartym na bramkach NOR będzie to realizowane poprzez przekształcenie funkcji XOR na postać przyjazną dla bramek NOR.

Generowanie przeniesienia na najmłodszym poziomie wynosi:
\begin{equation}
P_0 = A_0 \cdot B_0
\end{equation}

Dla bitów starszych:
\begin{equation}
S_1 = A_1 \oplus B_1 \oplus P_0
\end{equation}
oraz:
\begin{equation}
P_1 = (A_1 \cdot B_1) + (P_0 \cdot (A_1 \oplus B_1))
\end{equation}

Wynik najstarszy to:
\begin{equation}
S_2 = P_1
\end{equation}

Funkcje te można zapisać w postaci sumy iloczynów lub iloczynu sum, a następnie zminimalizować do postaci przyjaznej dla bramek NOR.

W praktyce realizacji na bramkach NOR każda z tych funkcji może być przekształcana do postaci, w której wykorzystywane są wyłącznie operacje NOR, co wymaga odpowiedniego zaznaczenia negacji i zamiany bramek na ich równoważniki.

Schemat działania można opisać następująco:
\begin{enumerate}
    \item poziom młodszy: określenie \(S_0\) oraz \(P_0\),
    \item poziom starszy: określenie \(S_1\) oraz \(P_1\),
    \item wyjście najstarsze: \(S_2 = P_1\).
\end{enumerate}

\subsection{Prezentacja układu i przejście do dalszej części pracy}

Struktura układu przedstawiona jest na schemacie blokowym oraz schematach logicznych. Na rysunku \ref{fig:schemat_bloku2} przedstawiono układ w formie uproszczonej, z wydzieleniem dwóch poziomów sumacyjnych oraz wyjścia przeniesienia. Na rysunku \ref{fig:schemat_logical2} szczegółowo przedstawiono zależności między wejściami a wyjściami w formie bramek logicznych.

\begin{figure}[h]
\centering
\begin{subfigure}[b]{0.8\textwidth}
    \centering
    \includegraphics[width=\textwidth]{syg1.png}
    \caption{Schemath blokowy układu}
    \label{fig:schemat_bloku2}
\end{subfigure}
\begin{subfigure}[b]{0.8\textwidth}
    \centering
    \includegraphics[width=\textwidth]{syg2.png}
    \caption{Schemath logiczny układu}
    \label{fig:schemat_logical2}
\end{subfigure}
\caption{Prezentacja układu sumatora 2-bitowego z 3-bitową sumą}
\label{fig:schematy_point2}
\end{figure}

Warto zwrócić uwagę na to, że w schemacie logicznym wejścia układu oraz wyjścia są połączone z bramkami NOR. Rysunek \ref{fig:schemat_bloku2} przedstawia układ w formie hierarchicznej: wejścia \(A_1, A_0, B_1, B_0\) są przetwarzane na poziomie młodszym, a następnie na poziomie starszym, a wynik końcowy jest sumą 3-bitową.

Następnie, w kolejnym punkcie pracy omówione zostaną minimalizacja funkcji logicznych na bramkach NOR, tablica prawdy oraz dalsze schematy i rysunki. Dalsza część pracy odnosi się bezpośrednio do przedstawionej tu struktury układu.

\begin{figure}[h]
\centering
\includegraphics[width=0.8\textwidth]{syg3.png}
\caption{Rysunek prezentujący przejście do następnego punktu pracy}
\label{fig:przejscie_point2}
\end{figure}

\end{document}
